\chapter{Introduction}
\ChapterLabel{Introduction}
\index{introduction}
\index{analog-to-digital conversion}
\index{source of truth}

\section{Chapter Responsibility}
\label{Section:IntroductionResponsibility}
Establish the purpose, scope, audience, and conventions for the embedded
analog-to-digital conversion system engineering record.

\noindent\textbf{Expected result:} A reviewable Week 1 scope that distinguishes
confirmed intent from decisions still to be made.

\section{Purpose}
\ProjectName\ is Jack Marino's embedded systems course project. Its confirmed
purpose is to accept an analog input and produce sampled numeric readings
using an \gls{adc}. Hardware has not yet been selected.
The architecture book records the project intent, system context, development
plan, and the decisions needed before detailed requirements and implementation.

The source evidence for this draft is the project owner's description and
follow-up confirming ADC readings and hardware TBD on November 10, 2026.
No previous implementation or measurement evidence was supplied for this
submission. The draft does not claim that hardware or firmware already exists.

\section{Scope}
The proposed initial demonstration covers the path from an analog input at
the system boundary through acquisition and conversion to an observable digital
reading at the output boundary. The particular input source, electrical limits,
channel count, conversion performance, and output method remain open in
\hyperref[Table:WeekOneTBDs]{the Week 1 decision register}.
The hardware--software boundary is described in \ChapterRef{SystemContext}.

The Week 1 deliverable is the Introduction, Development Plan, and System Context,
included in the complete architecture book. Later chapters retain the course
template to support subsequent work. Their service-request examples, example
requirements, and example tests are instructional material, not ADC-system
requirements, architecture, implemented functions, or evidence. Project-specific
requirements will be baselined after the Week 1 decisions are reviewed.

\begin{dsnNote}[Proposed scope limits for reviewer confirmation]
The initial demonstration is proposed as a supervised bench exercise.
Custom board manufacture, closed-loop actuation, production deployment,
remote connectivity, and certification are outside the proposed initial scope.
These are proposed limits, not approved exclusions. Safety, security, and
environmental applicability still require the review recorded in
\hyperref[Table:WeekOneTBDs]{the decision register}.
\end{dsnNote}

\section{Audience}
Jack Marino is the project owner and student author, responsible for resolving
project details, implementation, and evidence collection. The course reviewer
assesses scope, feasibility, document quality, and the proposed development
gates. A demonstration observer can later assess whether the digital readings
meet the agreed intended use. Any additional stakeholder remains TBD.

Readers can use this chapter for intent and document status,
\ChapterRef{SystemContext} for the boundary and external interactions, and
\ChapterRef{DevelopmentPlan} for ownership, milestones, risks, and decisions.

\section{Normative Language and Priority}
Binding requirements use \textbf{shall} in \ChapterRef{Requirements}.
Requirement priority uses \gls{must}, \gls{should}, \gls{could}, or \gls{would}
as defined by the template. Statements labeled proposed, TBD, or planned are
not approved requirements or evidence of completion. Week 1 introduces no
numerical acquisition requirement because no numerical target has been supplied.

\section{Document as Source of Truth}
The authored LaTeX and editable TikZ sources are the maintained engineering
record. The generated PDF is a review artifact built from those sources.
An implementation and its test evidence must identify the source baseline,
hardware revision, firmware revision, and test configuration before a later
claim of conformance. A successful document build establishes document
buildability; it does not establish ADC accuracy, firmware correctness, or
stakeholder acceptance.

\section{Initializing a New Project from This Template}
\label{Section:InitializeNewProject}
The existing \texttt{dsn} file namespace and required chapter filenames are
retained for this course submission. Project metadata identifies the ADC
project and its draft status. Unresolved information is centralized in
\hyperref[Table:WeekOneTBDs]{the Week 1 decision register}, rather than replaced
with invented hardware selections or performance values.
Further initialization follows
\hyperref[Skill:AIS-02]{AIS-02: Initialize a project}.

\section{Template Tailoring and Domain Profiles}
\label{Section:DomainProfiles}
The Embedded Systems profile is Applicable because the owner explicitly
identified an embedded ADC system. Its authoritative decision and rationale
are in \hyperref[Table:ProjectDomainProfiles]{the profile decision table}.
\ChapterRef{EmbeddedSystemsProfile} supplies hardware, timing, resource,
power, fault, safety, security, and lifecycle questions for later requirements
work. No extra domain profile is selected without evidence and human review.

\begin{table}[htbp]
\centering
\caption{Domain Profile Applicability}
\label{Table:DomainProfileApplicability}
\begin{tabular}{p{3.5cm}p{8cm}}
\toprule
Profile & Authoritative decision\\
\midrule
Embedded Systems & Applicable; see
\hyperref[Table:ProjectDomainProfiles]{Development Plan}.\\
\bottomrule
\end{tabular}
\end{table}

\section{Using This Template with AI Agents}
\label{Section:UsingTemplateAI}
AI assistance can draft and check this record using
\ChapterRef{AIAgentGuide} and \ChapterRef{AIImplementationWorkflow}.
The human owner confirms project facts and decisions; AI output does not
approve a baseline. Material inputs and unresolved assumptions are recorded in
\ChapterRef{AIPromptLog}. For this Week 1 change, one agent performs
coordination, drafting, build preparation, and a self-review. Human review
remains pending.

\section{Development and Evidence}
The project proceeds from reviewed context to measurable requirements, then
hardware and software design, implementation, and tests. Test findings feed
back into affected decisions and traceability before the next baseline.
The proposed milestones and evidence gates are in
\hyperref[Table:WeekOneMilestones]{the development milestones}. No gate is
reported passed without its required evidence and authorized disposition.

\begin{figure}[htbp]
\centering
\begin{tikzpicture}[>=Stealth,
 box/.style={draw=dsnBlue,fill=dsnLightBlue,align=center,text width=3.2cm,
 minimum height=0.85cm,font=\small}]
\node[box] (needs) at (-3,2) {Stakeholder needs};
\node[box] (req) at (-3,0) {Requirements};
\node[box] (design) at (-3,-2) {Design};
\node[box] (impl) at (0,-4) {Implementation};
\node[box] (test) at (3,-2) {Unit / integration tests};
\node[box] (verify) at (3,0) {System verification};
\node[box] (valid) at (3,2) {Stakeholder validation};
\draw[->] (needs.south)--(req.north);
\draw[->] (req.south)--(design.north);
\draw[->] (design.south)|-(impl.west);
\draw[->] (impl.east)-|(test.south);
\draw[->] (test.north)--(verify.south);
\draw[->] (verify.north)--(valid.south);
\draw[dashed,<->] (needs.east)--(valid.west);
\draw[dashed,<->] (req.east)--(verify.west);
\draw[dashed,<->] (design.east)--(test.west);
\end{tikzpicture}

\caption{Generic development-to-evidence relationships; the ADC project gates
are defined in the Development Plan.}
\label{Figure:VModel}
\end{figure}
Figure~\ref{Figure:VModel} illustrates the retained template's evidence
organization. It is a lifecycle illustration, not a completed project baseline.
